\chapter{Lezione 4 -- Il tempo nei sistemi distribuiti (II): tempo logico}

\section{Di cosa parliamo}

La Lezione 3 si chiudeva con un limite: gli orologi fisici, per quanto sincronizzati, non bastano a ordinare gli eventi. Qui si cambia approccio. Spesso non interessa \textbf{quando} è successo qualcosa, ma \textbf{in che ordine} e \textbf{cosa ha causato cosa}. Per questo bastano contatori, senza nessun orologio vero.

Percorso della lezione:

\begin{enumerate}
\def\labelenumi{\arabic{enumi}.}
\tightlist
\item
  la relazione \textbf{happened-before}, che formalizza ``questo evento può aver causato quello'';
\item
  il \textbf{clock logico scalare} (di Lamport), un numero per evento;
\item
  il \textbf{vector clock}, un vettore per evento, che cattura la causalità in modo esatto;
\item
  un'applicazione: la \textbf{mutua esclusione di Lamport}.
\end{enumerate}

Slide usate: versione A.A. 2024/25 (``Time II'').

\section{Perché serve il tempo logico}

\subsection{L'esempio della chat}

\begin{inbreve}

Se i messaggi arrivano nell'ordine sbagliato, una risposta può arrivare prima della domanda e il significato cambia. Serve un modo per sapere che un messaggio ``dipende'' da un altro.

\end{inbreve}

Scenario: una chat distribuita per un ambiente critico (nessun server centrale, niente copie dei messaggi su un server, deve funzionare anche se la rete perde la sincronia, quindi \textbf{modello asincrono}: niente orologi fisici affidabili).

Ordine reale:

\begin{enumerate}
\def\labelenumi{\arabic{enumi}.}
\tightlist
\item
  Alice: ``Se il prezzo di Apple scende, che facciamo?''
\item
  Mike: ``Vendiamo Apple e compriamo Amazon.''
\item
  Alice: ``Chi lo fa?''
\item
  Il capo di Bob: ``Lo fa Bob.''
\end{enumerate}

Se a Bob arrivano solo i messaggi 2 e 4, oppure il 3 dopo il 4, legge ``Vendiamo Apple'' seguito da ``Lo fa Bob'' e agisce \textbf{subito}, anche se la condizione (``se il prezzo scende'') non l'ha mai vista.

Cosa serve: quando Bob riceve il messaggio di Mike, deve sapere che c'è qualcosa che l'ha \textbf{causato} (la domanda di Alice) e che lui non ha ancora visto.

\textbf{Obiettivo della lezione}: dare timestamp agli eventi in modo che rispettino l'idea intuitiva di causalità.

\section{La relazione happened-before}

\begin{inbreve}

``\(e\) è successo prima di \(e'\)'' (\(e \to e'\)) vuol dire che \(e\) \textbf{può aver influenzato} \(e'\): o sono sullo stesso processo uno dopo l'altro, o c'è una catena di messaggi che porta da \(e\) a \(e'\). Se non c'è nessuna catena, i due eventi sono \textbf{concorrenti}: nessuno dei due può sapere dell'altro.

\end{inbreve}

\subsection{Definizione (Lamport)}

\(e \to e'\) se vale una di queste regole:

\begin{enumerate}
\def\labelenumi{\arabic{enumi}.}
\tightlist
\item
  \textbf{Ordine locale}: \(e\) ed \(e'\) sono sullo stesso processo \(p_i\) ed \(e\) viene prima (\(e \to_i e'\)).
\item
  \textbf{Invio-ricezione}: per ogni messaggio \(m\), \(\text{send}(m) \to \text{receive}(m)\).
\item
  \textbf{Transitività}: se \(e \to e''\) ed \(e'' \to e'\), allora \(e \to e'\).
\end{enumerate}

In pratica: \(e \to e'\) se nel diagramma spazio-tempo si può andare da \(e\) a \(e'\) \textbf{seguendo le linee dei processi in avanti e le frecce dei messaggi}.

\subsection{Eventi concorrenti}

Se né \(e \to e'\) né \(e' \to e\), gli eventi sono \textbf{concorrenti}: \(e \parallel e'\).

Per due eventi qualsiasi vale sempre \textbf{esattamente una} di tre cose: \(e \to e'\), oppure \(e' \to e\), oppure \(e \parallel e'\).

\textbf{Ordine parziale}: happened-before è un ordine \textbf{parziale} (poset), non totale: ci sono coppie di eventi che non sono ordinate (quelle concorrenti). Di conseguenza le sequenze di eventi che rispettano la causalità possono essere \textbf{più d'una} (si possono scambiare tra loro gli eventi concorrenti).

\section{Clock logico scalare (Lamport)}

\begin{inbreve}

Ogni processo tiene un contatore. Lo aumenta a ogni evento e lo allega ai messaggi; chi riceve un messaggio porta il suo contatore almeno al valore ricevuto. Così, se un evento causa un altro, il secondo ha sempre un numero più grande.

\end{inbreve}

Un \textbf{clock logico} è un contatore software che cresce sempre e non ha niente a che fare con l'ora vera. Ogni processo \(p_i\) ha il suo \(L_i\) e lo usa per dare un timestamp \(L(e)\) agli eventi.

\subsection{Le regole}

\begin{itemize}
\tightlist
\item
  \textbf{Inizio}: \(L_i = 0\).
\item
  \textbf{Evento locale o invio di} \(m\): \(L_i = L_i + 1\); se è un invio, \(m\) porta con sé il timestamp \(t = L_i\).
\item
  \textbf{Ricezione di} \(m\) \textbf{con timestamp} \(t\): \(L_i = \max(t, L_i) + 1\).
\end{itemize}

Esempio minimo: \(p_0\) manda un messaggio al suo primo evento (\(L = 1\)). \(p_1\) ha già fatto un evento (\(L = 1\)) e riceve: \(\max(1, 1) + 1 = 2\).

\subsection{Proprietà garantita}

\[
e \to e' \Rightarrow L(e) < L(e')
\]

Perché: lungo ogni passo di una catena causale il contatore cresce di almeno 1. Sullo stesso processo per la regola dell'incremento; lungo un messaggio perché il ricevente prende il massimo e poi aggiunge 1.

\subsection{Il limite}

Il viceversa \textbf{non vale}:

\[
L(e) < L(e') \not\Rightarrow e \to e'
\]

Esempio dalle slide: \(e^1_2\) e \(e^1_3\) sono concorrenti ma hanno timestamp diversi (2 e 1); \(e^1_1\) e \(e^1_3\) sono concorrenti e hanno lo stesso timestamp (1). Quindi, guardando solo due numeri, \textbf{non si capisce se due eventi sono concorrenti o causalmente legati}. Se \(L(e) < L(e')\), sappiamo solo che \(e'\) \textbf{non} può essere successo prima di \(e\).

\subsection{Ordine totale con gli ID}

A volte serve un ordine \textbf{totale} (tutti d'accordo su chi viene prima, anche per eventi concorrenti). Si usa la coppia (timestamp, ID del processo):

\[
(L, i) < (L', j) \iff L < L' \ \lor\ (L = L' \land i < j)
\]

A parità di timestamp vince l'ID più piccolo. È arbitrario per gli eventi concorrenti, ma rispetta la causalità e tutti lo calcolano allo stesso modo. Si usa nella mutua esclusione più sotto.

\section{Vector clock}

\subsection{L'idea: portarsi dietro la storia}

\begin{inbreve}

Invece di un solo numero, ogni evento si porta dietro ``quanti eventi conosco di ciascun processo''. Se il vettore di \(e'\) contiene tutto quello che sa il vettore di \(e\) (e qualcosa in più), allora \(e\) ha causato \(e'\), e vale anche il viceversa.

\end{inbreve}

Le slide ci arrivano per gradi:

\begin{enumerate}
\def\labelenumi{\arabic{enumi}.}
\tightlist
\item
  \textbf{Prima idea}: il timestamp di \(e\) è l'intera \textbf{storia} degli eventi che sono successi prima di \(e\) (incluso \(e\)), divisa per processo. All'invio aggiungo l'evento alla mia storia; alla ricezione prendo, per ogni processo, la versione più aggiornata della sua storia tra la mia e quella del messaggio.
\item
  \textbf{Confronto}: \(e \to e'\) se e solo se la storia di \(e'\) \textbf{contiene} quella di \(e\). Infatti se \(e\) compare nella storia di \(e'\), l'informazione su \(e\) è arrivata a \(e'\) attraverso una catena di messaggi, cioè proprio un cammino causale.
\item
  \textbf{Ottimizzazione}: la storia di ogni processo cresce sempre ``in coda'' (è sempre un prefisso di quella successiva). Quindi invece di tutta la lista basta la sua \textbf{lunghezza}. La storia diventa un vettore di \(n\) interi.
\end{enumerate}

\subsection{Definizione e regole}

Un \textbf{vector clock} per \(n\) processi è un array di \(n\) interi. \(p_i\) tiene \(V_i\):

\begin{itemize}
\tightlist
\item
  \(V_i[i]\) = quanti eventi ha fatto \(p_i\);
\item
  \(V_i[j]\) (con \(j \neq i\)) = quanti eventi di \(p_j\) conosce \(p_i\).
\end{itemize}

Regole:

\begin{itemize}
\tightlist
\item
  \textbf{Inizio}: \(V_i[j] = 0\) per ogni \(j\).
\item
  \textbf{Evento locale o invio}: \(V_i[i] = V_i[i] + 1\); l'invio allega al messaggio il vettore \(V_i\).
\item
  \textbf{Ricezione di} \(m\) \textbf{con vettore} \(V_t\): \(V_i[j] = \max(V_t[j], V_i[j])\) per ogni \(j\), poi \(V_i[i] = V_i[i] + 1\).
\end{itemize}

\subsection{Confrontare due vettori}

\begin{itemize}
\tightlist
\item
  \(V \le V'\) se \(V[j] \le V'[j]\) per \textbf{ogni} \(j\).
\item
  \(V < V'\) se \(V \le V'\) e \(V \neq V'\) (almeno una componente strettamente minore).
\item
  Se né \(V \le V'\) né \(V' \le V\) (ognuno ha una componente più grande dell'altro), i vettori sono \textbf{non confrontabili}.
\end{itemize}

La proprietà chiave, che lo scalare non aveva:

\[
e \to e' \iff V(e) < V(e')
\]

Quindi: \(V(e) < V(e')\) significa \(e \to e'\); vettori non confrontabili significano \(e \parallel e'\).

\textbf{Esempio} (dalle slide): \([\text{-},\text{-},1] < [4,3,2]\), quindi \(e^1_3 \to e^4_1\). Invece \([4,3,2]\) e \([3,\text{-},3]\) non sono confrontabili (4 \textgreater{} 3 ma 2 \textless{} 3), quindi \(e^4_1 \parallel e^3_3\). Il trattino ``-'' vale 0.

\subsection{Quale clock per cosa}

\begin{itemize}
\tightlist
\item
  \textbf{Scalare} → quando basta un ordine totale coerente con la causalità. Esempio: mutua esclusione di Lamport.
\item
  \textbf{Vettoriale} → quando bisogna sapere davvero se un evento dipende da un altro. Esempio: \emph{causal broadcast} (che risolve il problema della chat).
\end{itemize}

\section{Mutua esclusione di Lamport}

\subsection{Il problema}

\begin{inbreve}

C'è una risorsa (la \emph{sezione critica}, CS) che può usare un solo processo alla volta. I processi devono mettersi d'accordo, solo scambiandosi messaggi, su chi entra e in che ordine.

\end{inbreve}

L'astrazione ha tre eventi:

\begin{itemize}
\tightlist
\item
  \textbf{Request} (dall'alto): voglio entrare in CS.
\item
  \textbf{Grant} (verso l'alto): puoi entrare.
\item
  \textbf{Release} (dall'alto): sono uscito.
\end{itemize}

Proprietà:

\begin{itemize}
\tightlist
\item
  \textbf{Mutua esclusione} (safety): in ogni istante al massimo un processo è in CS.
\item
  \textbf{Liveness}: se un processo chiede di entrare, prima o poi entra.
\item
  \textbf{Fairness}: se la richiesta di \(p\) è \emph{happened-before} quella di \(q\), \(q\) non entra prima di \(p\).
\end{itemize}

Ipotesi dell'algoritmo: \textbf{nessun crash} (\(f = 0\)) e \textbf{link FIFO perfect} (dalla Lezione 2).

\subsection{L'idea}

\begin{inbreve}

Tutti tengono una coda delle richieste ordinata per (timestamp, ID). Entri quando la tua richiesta è in testa alla coda \textbf{e} hai sentito tutti gli altri dopo la tua richiesta: così sei sicuro che non esiste una richiesta più vecchia ancora in viaggio.

\end{inbreve}

Come funziona:

\begin{itemize}
\tightlist
\item
  \textbf{Richiesta}: \(p_i\) incrementa il clock, manda \(\text{REQ}\) con timestamp \((ck, i)\) a tutti e la mette nella sua coda \texttt{Requests}.
\item
  \textbf{Ricezione di una REQ da} \(p_j\): \(p_i\) la mette in \texttt{Requests} e risponde con un \(\text{ACK}\) con il suo timestamp attuale.
\item
  \textbf{Entrata in CS}, se valgono entrambe:

  \begin{enumerate}
  \def\labelenumi{\arabic{enumi}.}
  \tightlist
  \item
    la richiesta di \(p_i\) ha il timestamp \textbf{più piccolo} in \texttt{Requests};
  \item
    \(p_i\) ha ricevuto da \textbf{ogni} altro processo un messaggio con timestamp \textbf{maggiore} di quello della sua richiesta.
  \end{enumerate}
\item
  \textbf{Uscita}: \(p_i\) toglie la sua richiesta dalla coda e manda \(\text{RELEASE}\) a tutti.
\item
  \textbf{Ricezione di una RELEASE da} \(p_j\): \(p_i\) toglie la richiesta di \(p_j\) dalla coda.
\end{itemize}

Perché la condizione 2 + FIFO funziona: se \(p_j\) mi ha mandato un ACK con timestamp maggiore del mio, ogni richiesta che \(p_j\) ha fatto \textbf{prima} ha timestamp più piccolo e, per FIFO, mi è già arrivata. Ogni richiesta che fa \textbf{dopo} avrà timestamp più grande del mio. Quindi nessuna richiesta più vecchia della mia può ``spuntare'' in ritardo.

\subsection{Pseudocodice (ripulito)}

Nella slide i nomi \texttt{Requests}/\texttt{Request\_set} sono usati a caso e la richiesta propria non viene messa in coda; qui è sistemato secondo le regole della slide precedente.

\begin{verbatim}
upon event Init do
    Requests := ∅; Acks := ∅
    clock := 0; my_req := ⊥

upon event Request do                          // voglio la CS
    clock := clock + 1
    my_req := (REQ, ts = ⟨clock, i⟩)
    Requests := Requests ∪ {my_req}
    forall p_j ∈ Π do send FIFO(p_j, my_req)

upon event Release do                          // esco dalla CS
    Requests := Requests \ {my_req}; my_req := ⊥
    clock := clock + 1
    forall p_j ∈ Π do send FIFO(p_j, (RLS, ts = ⟨clock, i⟩))

// ts(x) < ts(y): clock più piccolo, a parità ID più piccolo
upon my_req ≠ ⊥
     and no req ∈ Requests with ts(req) < ts(my_req)
     and for every p ∈ Π exists m ∈ Acks from p with ts(m) > ts(my_req) do
    trigger Granted

upon event Deliver(m) do
    clock := max(clock(m), clock) + 1
    if m is REQ then
        Requests := Requests ∪ {m}
        clock := clock + 1
        send FIFO(sender(m), (ACK, ts = ⟨clock, i⟩))
    else if m is ACK then
        Acks := Acks ∪ {m}
    else if m is RLS then
        remove from Requests the request of sender(m)
\end{verbatim}

\subsection{Esempio (dalle slide, 3 processi)}

\begin{enumerate}
\def\labelenumi{\arabic{enumi}.}
\tightlist
\item
  \(p_3\) chiede la CS con \(ts = (3, 3)\). \(p_1\) e \(p_2\) la mettono in coda e rispondono con un ACK.
\item
  Poi anche \(p_1\) chiede, con \(ts = (6, 1)\). Riceve gli ACK di tutti, ma in coda c'è \((3, 3) < (6, 1)\): \textbf{aspetta}.
\item
  \(p_3\) ha tutti gli ACK ed è in testa alla coda: \textbf{entra}. Quando esce manda RELEASE.
\item
  Ricevuta la RELEASE, \(p_1\) toglie \((3,3)\) dalla coda: ora è in testa e ha gli ACK, quindi \textbf{entra}.
\end{enumerate}

\subsection{Perché funziona}

\textbf{Mutua esclusione} (per assurdo). Supponiamo che \(p_i\) e \(p_j\) siano in CS insieme. Entrambi hanno ricevuto l'ACK dell'altro ed entrambi hanno la propria richiesta in testa alla coda.

\begin{itemize}
\tightlist
\item
  Se \(p_j\) avesse fatto la sua richiesta \textbf{dopo} aver mandato l'ACK a \(p_i\), la richiesta di \(p_j\) avrebbe timestamp maggiore di quella di \(p_i\) (il clock di \(p_j\) è già passato oltre). Allora \(p_j\) vede in coda la richiesta di \(p_i\), più piccola, e non entra. Contraddizione.
\item
  Quindi \(p_j\) ha fatto la richiesta \textbf{prima} di mandare l'ACK. Per \textbf{FIFO} la REQ di \(p_j\) arriva a \(p_i\) prima dell'ACK, quindi è nella coda di \(p_i\). Simmetricamente la REQ di \(p_i\) è nella coda di \(p_j\).
\item
  Entrambi hanno entrambe le richieste in coda, e l'ordine (ts, ID) è totale: solo una delle due è la più piccola, quindi solo uno dei due può essere in testa. Contraddizione.
\end{itemize}

\textbf{Fairness.} Supponiamo che la REQ di \(p_j\) sia happened-before quella di \(p_i\), ma \(p_i\) entri prima. Allora \(ts(\text{REQ}_j) < ts(\text{REQ}_i)\) (proprietà del clock scalare). \(p_i\) entra solo dopo aver ricevuto l'ACK di \(p_j\), e per FIFO la REQ di \(p_j\) gli è già arrivata prima di quell'ACK. Quindi nella coda di \(p_i\) c'è una richiesta più piccola della sua, e non può entrare. Contraddizione.

\subsection{Costo}

\begin{itemize}
\tightlist
\item
  \textbf{Messaggi per ogni CS}: \(3(N-1)\), cioè \(N-1\) REQ, \(N-1\) ACK e \(N-1\) RELEASE.
\item
  \textbf{Ritardo nel caso migliore} (nessuno in CS, uno solo chiede): 2 tempi di messaggio, la REQ all'andata e l'ACK al ritorno.
\item
  \textbf{Limite}: se un solo processo crasha, gli altri aspettano per sempre il suo ACK. Non tollera nessun guasto.
\end{itemize}

\section{Esercizi e homework}

\begin{approfondimento}{Esercizio vector clock 1 (soluzione mia, controlla i calcoli)}

Vettori nell'ordine \([P_1, P_2, P_3, P_4]\), eventi da sinistra a destra.

\begin{itemize}
\tightlist
\item
  \(P_1\): invio a \(P_2\) \([1,0,0,0]\); ricezione da \(P_3\) \([2,0,1,0]\); invio a \(P_2\) \([3,0,1,0]\); \(e_3\) (ricezione da \(P_2\)) \([4,7,4,3]\).
\item
  \(P_2\): ric. \([1,1,0,0]\); invio a \(P_4\) \([1,2,0,0]\); ric. da \(P_1\) \([3,3,1,0]\); invio a \(P_3\) \([3,4,1,0]\); ric. da \(P_3\) \([3,5,4,1]\); ric. da \(P_4\) \([3,6,4,3]\); invio a \(P_1\) \([3,7,4,3]\).
\item
  \(P_3\): \(e_1\) (invio a \(P_1\)) \([0,0,1,0]\); ric. da \(P_4\) \([0,0,2,1]\); ric. da \(P_2\) \([3,4,3,1]\); invio a \(P_2\) \([3,4,4,1]\).
\item
  \(P_4\): invio a \(P_3\) \([0,0,0,1]\); ric. da \(P_2\) \([1,2,0,2]\); \(e_2\) (invio a \(P_2\)) \([1,2,0,3]\).
\end{itemize}

Confronti:

\begin{itemize}
\tightlist
\item
  \(e_1 \to e_2\): \([0,0,1,0]\) vs \([1,2,0,3]\). La terza componente è 1 \textgreater{} 0: non confrontabili, quindi \textbf{concorrenti}. \textbf{FALSA}.
\item
  \(e_1 \to e_3\): \([0,0,1,0] < [4,7,4,3]\). \textbf{Vera}.
\item
  \(e_2 \to e_3\): \([1,2,0,3] < [4,7,4,3]\). \textbf{Vera}.
\end{itemize}

\end{approfondimento}

\begin{approfondimento}{Esercizio vector clock 2 (soluzione mia, controlla i calcoli)}

\begin{itemize}
\tightlist
\item
  \(P_1\): invio a \(P_2\) \([1,0,0,0]\); \(e_1\) (invio a \(P_4\)) \([2,0,0,0]\); invio a \(P_2\) \([3,0,0,0]\); \(e_3\) (ric. da \(P_3\)) \([4,4,4,3]\).
\item
  \(P_2\): \(e_2\) (ric. da \(P_1\)) \([1,1,0,0]\); ric. da \(P_4\) \([1,2,0,1]\); ric. da \(P_3\) \([1,3,1,1]\); invio a \(P_3\) \([1,4,1,1]\); invio a \(P_4\) \([1,5,1,1]\); \(e_4\) (ric. da \(P_1\)) \([3,6,1,1]\).
\item
  \(P_3\): invio a \(P_2\) \([0,0,1,0]\); ric. da \(P_2\) \([1,4,2,1]\); ric. da \(P_4\) \([2,4,3,3]\); invio a \(P_1\) \([2,4,4,3]\).
\item
  \(P_4\): invio a \(P_2\) \([0,0,0,1]\); ric. \(e_1\) \([2,0,0,2]\); invio a \(P_3\) \([2,0,0,3]\); \(e_5\) (ric. da \(P_2\)) \([2,5,1,4]\).
\end{itemize}

Confronti:

\begin{itemize}
\tightlist
\item
  \(e_1 \to e_2\): \([2,0,0,0]\) vs \([1,1,0,0]\): non confrontabili. \textbf{FALSA} (\(e_1 \parallel e_2\)).
\item
  \(e_1 \to e_3\): \([2,0,0,0] < [4,4,4,3]\). \textbf{Vera}.
\item
  \(e_2 \to e_3\): \([1,1,0,0] < [4,4,4,3]\). \textbf{Vera}.
\item
  (versione italiana) \(e_1 \to e_4\): \([2,0,0,0] < [3,6,1,1]\). \textbf{Vera}. \(e_1 \to e_5\): \([2,0,0,0] < [2,5,1,4]\). \textbf{Vera}.
\end{itemize}

\end{approfondimento}

\begin{approfondimento}{Break time: debugging di p1 con i vector clock (soluzione mia)}

\textbf{Dati}: 4 processi, clock incrementato prima dell'invio. \(p_1\) invia un messaggio con \((1,0,0,0)\) e poi riceve un messaggio con \((1,2,3,1)\).

\textbf{Q3.1: un'esecuzione possibile.}

\begin{enumerate}
\def\labelenumi{\arabic{enumi}.}
\tightlist
\item
  \(p_1\) manda a \(p_2\): \((1,0,0,0)\).
\item
  \(p_4\) manda a \(p_3\): \((0,0,0,1)\).
\item
  \(p_3\) riceve da \(p_4\): \((0,0,1,1)\).
\item
  \(p_2\) riceve da \(p_1\): \((1,1,0,0)\), poi manda a \(p_3\): \((1,2,0,0)\).
\item
  \(p_3\) riceve da \(p_2\): \((1,2,2,1)\), poi manda a \(p_1\): \((1,2,3,1)\).
\end{enumerate}

\textbf{Non è unica}: ad esempio si possono scambiare l'ordine dei passi 2-3 e 4, oppure il messaggio di \(p_4\) può passare da \(p_2\) invece che arrivare direttamente a \(p_3\).

\textbf{Q3.2: un processo può non aver mandato nessun messaggio?} No.~L'informazione su un processo arriva agli altri \textbf{solo} dentro i messaggi. Ogni componente del vettore ricevuto è diversa da 0, quindi chi ha mandato il messaggio conosce eventi di tutti e quattro i processi.

\begin{itemize}
\tightlist
\item
  \(p_1\) ha mandato il suo messaggio iniziale.
\item
  Il mittente finale non può essere \(p_4\): con la componente \(p_4 = 1\) il suo unico evento sarebbe l'invio stesso, ma allora non potrebbe conoscere gli eventi degli altri (ricevere è un evento e farebbe salire la sua componente).
\item
  Quindi il mittente è \(p_2\) o \(p_3\). L'altro dei due e \(p_4\) hanno componente diversa da 0 nel vettore del mittente, quindi devono avergli fatto arrivare informazioni: hanno spedito almeno un messaggio. E il mittente ha spedito il messaggio finale. Tutti hanno spedito.
\end{itemize}

\end{approfondimento}

\begin{approfondimento}{Homework: vector clock senza conoscere i processi (soluzione proposta,
non nelle slide)}

\textbf{Idea}: al posto di un array di \(n\) posizioni si usa una \textbf{mappa} ID → contatore. Un ID assente vale 0.

\begin{itemize}
\tightlist
\item
  Evento locale o invio: \texttt{V{[}my\_id{]}\ :=\ V{[}my\_id{]}\ +\ 1}.
\item
  Ricezione di \(V_t\): per ogni ID presente in \(V\) o in \(V_t\), \texttt{V{[}id{]}\ :=\ max(V{[}id{]},\ V\_t{[}id{]})} (gli ID nuovi vengono aggiunti); poi \texttt{V{[}my\_id{]}\ :=\ V{[}my\_id{]}\ +\ 1}.
\item
  Confronto: come prima, ma sull'unione delle chiavi delle due mappe.
\end{itemize}

\textbf{Eventi concorrenti}: ciascuna delle due mappe ha almeno un ID con valore più alto dell'altra, spesso un ID che nell'altra \textbf{non compare proprio} (un processo di cui l'altro evento non ha mai sentito parlare).

\textbf{Due processi con lo stesso ID}: i loro eventi finiscono nello stesso contatore, come se fossero un solo processo. Si rompe il ``⇐'': due eventi concorrenti dei due processi possono sembrare ordinati (es. \(\{X:1\}\) e \(\{X:2\}\) sembrano causalmente legati anche se non lo sono). In più i contatori non contano più correttamente gli eventi di nessuno dei due.

\end{approfondimento}

\begin{approfondimento}{Homework: Lamport ME senza link FIFO (soluzione proposta, non nelle
slide)}

\textbf{Mutua esclusione: violata.} La dimostrazione usa FIFO per dire che la REQ di \(p_j\) arriva prima del suo ACK. Senza FIFO un ACK può superare la REQ.

Controesempio con 2 processi:

\begin{enumerate}
\def\labelenumi{\arabic{enumi}.}
\tightlist
\item
  \(p_1\) ha già fatto altri eventi (clock = 5) e chiede la CS: \(\text{REQ}_1\) con \(ts = (6, 1)\).
\item
  \(p_2\) (clock = 0) chiede la CS: \(\text{REQ}_2\) con \(ts = (1, 2)\), che verso \(p_1\) viaggia lentissima.
\item
  \(p_2\) riceve \(\text{REQ}_1\) e risponde con un ACK con timestamp 8. L'ACK \textbf{supera} \(\text{REQ}_2\) e arriva a \(p_1\) per primo.
\item
  \(p_1\): in coda ha solo la sua richiesta, e ha un ACK di \(p_2\) con ts 8 \textgreater{} 6. \textbf{Entra}.
\item
  Arriva \(\text{REQ}_2\); \(p_1\) risponde con un ACK (ts \textgreater{} 1). \(p_2\) ha in coda \((1,2)\) e \((6,1)\): la sua è la più piccola e ha l'ACK. \textbf{Entra anche lui}. Due processi in CS: violazione.
\end{enumerate}

\textbf{Fairness: può essere violata} per lo stesso motivo: la sua dimostrazione usa FIFO per garantire che la richiesta più vecchia sia già in coda quando arriva l'ACK.

\textbf{Liveness: sembra reggere.} I link sono comunque perfect, quindi tutti i messaggi arrivano. Una richiesta arrivata in ritardo può rimescolare l'ordine, ma prima o poi ogni processo riceve tutto e la richiesta più piccola rimasta può entrare. Da verificare a lezione.

\end{approfondimento}

\section{Riepilogo}

\begin{itemize}
\tightlist
\item
  Spesso conta l'\textbf{ordine causale}, non l'ora. \textbf{Happened-before} (\(\to\)): ordine locale + invio prima di ricezione + transitività. È un ordine \textbf{parziale}; gli eventi non ordinati sono \textbf{concorrenti}.
\item
  \textbf{Clock scalare} (Lamport): \(+1\) a ogni evento, \(\max + 1\) alla ricezione. Garantisce \(e \to e' \Rightarrow L(e) < L(e')\) ma \textbf{non} il viceversa. Con (ts, ID) si ottiene un ordine totale.
\item
  \textbf{Vector clock}: un contatore per processo, \(\max\) componente per componente alla ricezione. Garantisce \(e \to e' \iff V(e) < V(e')\); vettori non confrontabili = eventi concorrenti.
\item
  \textbf{Mutua esclusione di Lamport}: coda ordinata per (ts, ID); entri se sei in testa e hai sentito tutti dopo la tua richiesta. Serve FIFO, nessun crash tollerato, \(3(N-1)\) messaggi per CS.
\end{itemize}
